
% -------------------------------------------------------
% INSERT
% -------------------------------------------------------

\subsubsection{I2: Bulk import ocen}
\begin{figure}[H]
    \centering
    \includegraphics[width=0.8\textwidth]{charts/small/chart_I2.png}
    \caption{I2 -- Bulk import ocen}
\end{figure}

MongoDB jest tu wyraźnie najszybszy (36~ms vs 190--290~ms w pozostałych
silnikach) -- wstawianie dokumentów do kolekcji bez walidacji schematu
jest tańsze niż klasyczny \texttt{INSERT} z kontrolą kluczy obcych
i typów. Założenie indeksów spowalnia Mongo \emph{trzykrotnie}
(z 36 do 111~ms), ponieważ każdy wstawiany dokument musi
zaktualizować strukturę indeksu -- co dobrze potwierdza drugą część
hipotezy pracy. PostgreSQL i MySQL prawie nie odczuwają tej różnicy
między wariantami, bo ich nadrzędna prędkość \texttt{INSERT}-a
i tak jest ograniczona przez WAL/redo log, a nie przez liczbę
indeksów dodatkowych.

\subsubsection{I3: Dodanie serialu z pełnym drzewem}
\begin{figure}[H]
    \centering
    \includegraphics[width=0.8\textwidth]{charts/small/chart_I3.png}
    \caption{I3 -- Dodanie serialu z pełnym drzewem}
\end{figure}

Tu MongoDB \emph{miażdży} konkurencję -- 0.9~ms wobec 17--25~ms
w pozostałych bazach. Powód jest prosty: w modelu dokumentowym serial
z sezonami i odcinkami to \emph{jeden} dokument zapisywany jednym
wywołaniem \texttt{insert\_one}. W bazach relacyjnych ten sam serial
wymaga osobnych \texttt{INSERT}-ów do tabel \texttt{content},
\texttt{seasons} i \texttt{episodes}, a w Neo4j -- wielu
\texttt{CREATE} dla węzłów i relacji. Indeksy nie wpływają tu
istotnie -- wszystkie operacje są \texttt{INSERT}-ami pojedynczych
encji rzędu kilkunastu rekordów, a koszt zapisu indeksu na takiej
wielkości jest w granicach szumu pomiarowego.

% -------------------------------------------------------
% SELECT
% -------------------------------------------------------

\subsubsection{S2: Filtrowanie kolaboratywne}
\begin{figure}[H]
    \centering
    \includegraphics[width=0.8\textwidth]{charts/small/chart_S2.png}
    \caption{S2 -- Filtrowanie kolaboratywne}
\end{figure}

MySQL wygrywa we wszystkich wariantach -- dzięki temu, że indeks na
\texttt{profile\_id} jest niemożliwy do upuszczenia (chroni FK -- patrz
uwaga w 11.1), MySQL i tak korzysta z optymalnego planu nawet w
wariancie ,,bez indeksów''. Najmocniejszy zysk daje MongoDB,
który spada ze 158~ms do 16~ms -- to potwierdza obserwację z 11.3,
że bez indeksu na \texttt{profile\_id} Mongo musi skanować całą
kolekcję \texttt{watch\_history}. PostgreSQL halnie czas (z 78 do
40~ms), ale dalej zostaje mu rezydualny \texttt{Seq Scan} z
podzapytania \texttt{NOT IN}. Neo4j najmniej zyskuje, bo i tak
operuje na strukturze grafu, w której indeksy na atrybutach mają
mniejsze znaczenie niż przejścia po krawędziach.

\subsubsection{S3: TOP 100 według liczby ocen w okresie}
\begin{figure}[H]
    \centering
    \includegraphics[width=0.8\textwidth]{charts/small/chart_S3.png}
    \caption{S3 -- TOP 100 według liczby ocen w okresie}
\end{figure}

PostgreSQL prowadzi (14--16~ms), za nim MongoDB (26--28~ms),
a najwolniejsze są MySQL i Neo4j (38--71~ms). Kluczowa obserwacja:
\textbf{indeks praktycznie nie pomaga} -- w żadnym z silników nie ma
zauważalnej różnicy między wariantami, w MySQL i PostgreSQL nawet
występuje minimalne pogorszenie. Zgadza się to z analizą EXPLAIN
z 11.3: filtr po dacie zwraca tak szeroki zbiór wierszy, że indeks
na \texttt{started\_at} nie zmienia faktu, że trzeba przeczytać
większość tabeli. Ten scenariusz pokazuje granicę użyteczności
indeksu -- dla zapytań z niską selektywnością strukturalna optymalizacja
zapytania (np. agregacje) jest istotniejsza niż obecność indeksu.

\subsubsection{S4: Wyszukiwanie pełnotekstowe po tytule}
\begin{figure}[H]
    \centering
    \includegraphics[width=0.8\textwidth]{charts/small/chart_S4.png}
    \caption{S4 -- Wyszukiwanie pełnotekstowe po tytule}
\end{figure}

Na małym wolumenie wszystkie cztery silniki radzą sobie szybko
(1.4--7~ms), więc różnice między wariantami są mało spektakularne.
Najwolniejsze jest Neo4j (6.7~ms) -- skan wszystkich węzłów
\texttt{:Content} jest w tej skali odczuwalny. PostgreSQL i MySQL
bez indeksu kończą w niespełna 2~ms, bo w tabeli jest tylko
2\,000 tytułów do przeszukania. Prawdziwa wartość indeksów
pełnotekstowych ujawni się dopiero na wolumenach \emph{medium}
i \emph{large}, gdzie różnica między skanem pełnym a indeksem
sięga rzędów wielkości.

% -------------------------------------------------------
% DELETE
% -------------------------------------------------------

\subsubsection{D1: Usunięcie treści z kaskadą}
\begin{figure}[H]
    \centering
    \includegraphics[width=0.8\textwidth]{charts/small/chart_D1.png}
    \caption{D1 -- Usunięcie treści z kaskadą}
\end{figure}

PostgreSQL pokazuje tu drugą połowę hipotezy w dobrym świetle:
bez indeksów 17~ms, z indeksami \emph{1.2~ms} -- przyspieszenie
o rząd wielkości. Powodem jest to, że kaskadowe usunięcie
(\texttt{ON DELETE CASCADE}) wymaga znalezienia wszystkich wierszy
potomnych po kluczu obcym -- bez indeksu na FK to oznacza pełen
skan tabel \texttt{watch\_history}, \texttt{ratings} i \texttt{my\_list}.
MongoDB jest tu najwolniejszy w obu wariantach (91~ms / 22~ms),
ponieważ usunięcie treści wymaga skasowania wszystkich powiązanych
dokumentów w innych kolekcjach. MySQL ma stały czas 6~ms, bo indeks
na FK jest u niego zawsze obecny. Neo4j 11--12~ms -- relacje wskazują
na właściwe węzły, więc kaskada to po prostu usunięcie tych krawędzi.

\subsubsection{D4: Usunięcie pozycji z listy}
\begin{figure}[H]
    \centering
    \includegraphics[width=0.8\textwidth]{charts/small/chart_D4.png}
    \caption{D4 -- Usunięcie pozycji z listy}
\end{figure}

Najprostsza operacja w zestawie -- usunięcie jednego wiersza po
kluczu głównym. PostgreSQL kończy w 1~ms w obu wariantach,
MongoDB z indeksem zjeżdża do 0.6~ms. Bez indeksu MongoDB potrzebuje
20~ms -- klasyczny przypadek, w którym \texttt{deleteOne} bez
indeksu i tak musi przeskanować kolekcję, żeby znaleźć dokument
do usunięcia. To samo zjawisko będzie się skalować liniowo
z wolumenem (zobaczymy 200~ms na \emph{large}). MySQL i Neo4j
działają stabilnie 5--8~ms.

\subsubsection{D6: Masowe usunięcie nieaktywnych użytkowników}
\begin{figure}[H]
    \centering
    \includegraphics[width=0.8\textwidth]{charts/small/chart_D6.png}
    \caption{D6 -- Masowe usunięcie nieaktywnych użytkowników}
\end{figure}

Pierwszy scenariusz, w którym brak indeksu boli naprawdę mocno --
PostgreSQL bez indeksu na \texttt{users.status} kończy w 37~ms,
a z indeksem w 1.2~ms (30-krotnie szybciej). To kluczowy wynik
wspierający hipotezę: indeks dramatycznie przyspiesza znajdowanie
podzbioru wierszy do usunięcia, bo zamiast pełnego skanu \texttt{users}
silnik trafia od razu na rekordy ze statusem ,,inactive''. MongoDB
i Neo4j również zyskują (0.8 i 7~ms odpowiednio), ale dystans nie
jest tak spektakularny. \emph{Spektakularne} skutki braku indeksu
zobaczymy dopiero na większych wolumenach -- na \emph{large}
PostgreSQL bez indeksu osiągnie pomiar mierzony w \emph{sekundach}.

% -------------------------------------------------------
% UPDATE
% -------------------------------------------------------

\subsubsection{U2: Przeliczenie średniej oceny}
\begin{figure}[H]
    \centering
    \includegraphics[width=0.8\textwidth]{charts/small/chart_U2.png}
    \caption{U2 -- Przeliczenie średniej oceny}
\end{figure}

PostgreSQL z indeksem schodzi z 3.7~ms do 1.3~ms, MongoDB z 11~ms
do 1.7~ms. Operacja przelicza średnią ocenę wybranego filmu
i wykonuje \texttt{UPDATE} pojedynczego wiersza/dokumentu --
indeks pomaga zlokalizować zarówno źródłowe oceny (\texttt{ratings})
jak i wiersz w \texttt{content}. MySQL ma stały czas 6~ms, Neo4j
4--6~ms. To typowy operacyjny scenariusz, w którym indeks działa
zgodnie z oczekiwaniami -- punktowy filtr + pojedyncza modyfikacja
to klasyczny ,,słodki punkt'' indeksu dodatkowego, czyli wąski filtr
\texttt{WHERE} bez konfliktu z dużymi zapisami.

\subsubsection{U3: Masowa zmiana planu subskrypcji}
\begin{figure}[H]
    \centering
    \includegraphics[width=0.8\textwidth]{charts/small/chart_U3.png}
    \caption{U3 -- Masowa zmiana planu subskrypcji}
\end{figure}

Najciekawszy wynik tej kategorii i \emph{dowód kosztu indeksów dla
operacji zapisu}. PostgreSQL bez indeksu kończy w 8.9~ms, z indeksem
w \emph{16.9~ms} -- czyli założenie indeksów pogorszyło wydajność
prawie dwukrotnie. To samo widać w MySQL (38 → 52~ms), Mongo
(29 → 33~ms) i Neo4j (53 → 43~ms tylko z lekkim odchyleniem).
Operacja modyfikuje pole \texttt{plan\_name} w wielu wierszach
\texttt{subscriptions} -- każdy zaktualizowany wiersz wymaga
przepisania w indeksie \texttt{idx\_subscriptions\_status\_plan},
co dodaje stały narzut do każdej zmiany. Ten scenariusz najwyraźniej
ze wszystkich potwierdza drugą część hipotezy: indeksy podnoszą
koszt operacji \texttt{UPDATE} na kolumnach indeksowanych.
